%%%%%%%%%%%%%%%%%%%%%%%%%%%%%%%%%%%%%%%%%%%%%%%%%%%%%%%%%%%%%%%%%%%%%%%%%%%%%%%%
%%  BACKGROUND AND RELATED WORK
%%%%%%%%%%%%%%%%%%%%%%%%%%%%%%%%%%%%%%%%%%%%%%%%%%%%%%%%%%%%%%%%%%%%%%%%%%%%%%%%
\chapter{Background and Related Work}\label{background}
The history of robotics traces a gradual shift from rigid, pre-scripted automation toward systems capable of flexible interaction with unstructured settings and human collaborators. As Li \emph{et al.}~\cite{Li2020} note, traditional manufacturing robots were high-precision, deterministic tools within strict safety zones, where human-robot interaction was kept to a minimum to avoid injury. In these environments, robots executed pre-programmed trajectories with little to no capacity for environmental sensing or real-time adjustment. This baseline lays the groundwork for the subsequent move toward interactive systems.

\gls{hrc} research sought to overcome these constraints and enable robots to work alongside humans. Even so, early systems struggled with the semantic gap, the fundamental difficulty of translating high-level human intent into low-level motor commands that a robot can reliably execute.

The response to this gap was \gls{eai}, which grounds intelligence in a physical embodiment and requires constant interaction with the environment \cite{Li2026}. In this view, an embodied agent must perceive, act upon, and learn from the physical world rather than reason over abstract symbols alone.

But learning from the physical world at scale is exactly where this ambition meets a hard constraint, as diverse, real-world interaction data is slow and costly to collect. This bottleneck spurred the development of large-scale procedural generation frameworks such as ProcTHOR~\cite{deitke2022procthorlargescaleembodiedai}. These systems train agents across thousands of auto-generated indoor environments. More recent work, such as the Unity Reinforcement Learning Playground~\cite{ye2025unityrlplaygroundversatile}, further refines simulation-based training by introducing versatile \gls{rl} frameworks for mobile robots and explicit \gls{s2r} support. Together, these approaches treat simulation not as a simplified proxy but as a structured pipeline toward real deployment.

The use of \glspl{fm}, which extend language models to encompass diverse sensory modalities such as vision and audio, has proven essential for the current generation of \gls{hrc} systems. Building on this integration, the remainder of this chapter traces the field's main phases, from semantic reasoning and grasp planning to end-to-end paradigms, edge deployment limits, skill-library composition, structured execution, hierarchical safety, fleet-scale management, and multi-robot coordination.

% ------------------------------------------------------------------------------
\section{Semantic Reasoning and Task Planning}
Early demonstrations such as ChatGPT for Robotics by Vemprala \emph{et al.}~\cite{vemprala2023chatgptroboticsdesignprinciples} and Voyager by Wang \emph{et al.}~\cite{wang2023voyageropenendedembodiedagent} established that \glspl{llm} could serve as high-level task planners, generating code sequences or natural language plans to split complex tasks. These works revealed the potential of language as an interface between human intent and robotic execution. At the same time, they exposed the limitations of text-only reasoning in physically grounded settings.

Liang \emph{et al.}~\cite{liang2023codepolicieslanguagemodel} addressed the translation of plans into executable policies in \gls{cap}, demonstrating that code-writing \glspl{llm} can translate natural language commands into robot policy code that chains perception APIs with control primitives. Unlike plan-based approaches that assume a fixed skill library, \gls{cap} writes parameterized control logic on the fly, thereby generalizing to unseen instructions without task-specific training. The key insight is that representing code as policies rather than as sequences of predefined skills enables the system to compose novel behaviors at runtime, advancing planning into execution.

Where \gls{cap} operates at the level of individual workspaces, Rana \emph{et al.}~\cite{rana2023sayplangroundinglargelanguage} tackled the challenge of scaling language-based planning to large, multi-room environments with SayPlan. By grounding an \gls{llm} in 3D scene graphs, SayPlan enables task planning within environments that are far too large for a single model to process in its entirety. Semantic search over a collapsed scene-graph representation expands only the task-relevant subgraph, keeping the model’s context window manageable while preserving awareness of the global structure; a separate iterative replanning loop then validates each draft plan against a scene-graph simulator and feeds infeasible actions back for correction. This is relevant for multi-robot coordination, where agents must reason about objects and workspaces beyond their immediate field of view, extending the same planning logic to larger settings.

% ------------------------------------------------------------------------------
\section{Grasp Planning and Physical Grounding}
While semantic planners provide the strategic plan for complex tasks, effective execution relies on the robot’s ability to physically engage with discrete objects. High-level planners, therefore, define what must be done, whereas grasp planning determines the specific mechanics of physically engaging the environment. As the survey by Newbury \emph{et al.}~\cite{newbury2023deeplearningapproachesgrasp} points out, the field has moved from analytic force-closure methods toward data-driven approaches that learn grasp quality directly from sensor observations. Large-scale benchmarks such as GraspNet-1Billion~\cite{graspnet} and frameworks such as Contact-GraspNet~\cite{sundermeyer2021contactgraspnetefficient6dofgrasp} and AnyGrasp~\cite{fang2023anygrasprobustefficientgrasp} drove this shift, with the latter supporting real-time 7-DOF candidate generation for unseen objects in cluttered scenes.

Given this system’s infrastructure-independent deployment goals, inference speed and deployment simplicity are decisive considerations. \gls{vgn}~\cite{breyer2021volumetricgraspingnetworkrealtime} addresses these constraints by representing the scene as a \gls{tsdf} voxel grid. Using this grid, it predicts grasp quality, orientation, and gripper width in a single forward pass, reducing planning time to approximately 10 ms. This approach also provides implicit collision awareness and zero-shot transfer from simulation to real hardware.

% ------------------------------------------------------------------------------
\section{End-to-End Control}
Building on this progression, the field subsequently shifted toward \gls{vla} models, in which a single architecture consumes visual and linguistic inputs to directly produce motor commands. PaLM-E~\cite{driess2023palmeembodiedmultimodallanguage} injects images, state estimates, and other modalities directly into the embedding space of a pre-trained language model, enabling grounded, multi-embodiment robotic planning. A single model, scaled up to 562 billion parameters, was trained across multiple robotic manipulation domains, demonstrating positive transfer across the visual-language and embodied reasoning domains.

Extending this line of work, RT-2~\cite{brohan2023rt2visionlanguageactionmodelstransfer} showed that \gls{vla} models trained on internet-scale vision-language data can transfer broad semantic knowledge, including logical reasoning and common-sense understanding, to robot manipulation tasks. This enabled generalization to object categories and scenarios completely absent from the robot’s physical training data.

% ------------------------------------------------------------------------------
\section{Computing Constraints and Edge Deployment}
A major challenge in moving \gls{eai} beyond the laboratory is the tension between model scale and the latency constraints of physical systems. End-to-end models such as RT-2~\cite{brohan2023rt2visionlanguageactionmodelstransfer} and PaLM-E~\cite{driess2023palmeembodiedmultimodallanguage} achieve impressive semantic generalization, but their parameter counts require computational infrastructure incompatible with mobile or resource-limited platforms.

As Grover \emph{et al.}~\cite{grover2026embodiedfoundationmodelsedge} and Kim \emph{et al.}~\cite{kim2024openvlaopensourcevisionlanguageactionmodel} note, running a large \gls{vla} is slow in two ways. Interpreting camera and language input is computationally expensive, and producing each action one step at a time is constrained by memory speed. Together, these constraints keep most models to fewer than five updates per second. That is too slow for close collaboration.

Even edge-optimized models such as TinyVLA~\cite{wen2025tinyvlafastdataefficientvisionlanguageaction} face a persistent trade-off between reasoning accuracy and real-time reaction. This infrastructure bottleneck motivates the modular approach we take. By separating visual comprehension from physical execution, our system keeps control frequencies predictable and uses the model for semantic guidance only when a new task requires planning.

% ------------------------------------------------------------------------------
\section{Skill Libraries and Retrieval-Augmented Control}
As the computational and data demands of end-to-end \gls{vla} models became apparent, a parallel approach emerged. Rather than generating low-level control directly, \glspl{llm} are used to sequence predefined, physically verified skills from libraries. This concept adapts \gls{rag}~\cite{lewis2021retrievalaugmentedgenerationknowledgeintensivenlp} from natural language processing to the physical domain; instead of retrieving text documents to ground a model’s answers, the system retrieves parameterized control primitives or executable code blocks.

Early implementations relied on affordance models to constrain skill selection. SayCan~\cite{ahn2022icanisay} showed that a \gls{llm} can decompose high-level instructions into skill sequences, utilizing a pre-trained affordance function to evaluate whether each skill was physically feasible in the current state. Voyager~\cite{wang2023voyageropenendedembodiedagent} extended this concept, demonstrating that agents could not only retrieve executable skills from a library but also iteratively write, refine, and store new ones based on semantic similarity.

Mower \emph{et al.}~\cite{Mower_2026} carried this onto standard robotics middleware with \gls{ros}-\gls{llm}, exposing an atomic action library as documented \gls{ros} actions and services and mapping sensor readings to text, so that the full library and the current scene enter the model's prompt; non-experts extend the library by demonstration. Its open-source models and atomic-action decomposition make it the closest published precedent for the stack we build here, but its correction loop runs through a human operator who types the repair between steps. We close that loop autonomously, re-querying the parser with the failed operation's own recovery suggestions (Section~\ref{sec:reflection}). It also programs a single arm, leaving composition across two arms sharing a workspace outside its scope.

By separating high-level reasoning from low-level execution, these modular architectures aim to avoid the latency and instability of generating motor commands on the fly. The central challenge remains mapping visual scene context to the appropriate subset of atomic operations in unstructured settings. We address this directly with a retrieval-augmented operations pool grounded in structured scene representations.

% ------------------------------------------------------------------------------
\section{Structured Execution and Sim-to-Real Transfer}
Where the skill libraries above sequence hand-built primitives, a further line of work uses \glspl{llm} to automate the training pipeline behind reliable execution. Eureka~\cite{ma2024eurekahumanlevelrewarddesign} showed that \glspl{llm} can autonomously design reward functions for \gls{rl}, outperforming human-engineered rewards on 83\% of tasks across 29 environments and 10 robot morphologies without using task-specific prompt engineering. DrEureka~\cite{ma2024dreurekalanguagemodelguided} extended this approach to the full \gls{s2r} pipeline by utilizing an \gls{llm} to generate domain-randomization configurations, thereby automating two of the most labor-intensive manual steps in simulation-based policy training.

% ------------------------------------------------------------------------------
\section{Hierarchical Refinement and Safety}\label{sec:refinement_and_safety}
To address the hallucinations and safety risks that arise when models operate in physical environments, subsequent architectures introduced mechanisms for structured self-correction and transparent reasoning. Shinn \emph{et al.}~\cite{shinn2023reflexionlanguageagentsverbal} proposed Reflexion, in which an agent turns its own failures into natural-language feedback and conditions its next attempt on that critique rather than on gradient updates. This retry-on-reflection loop is the direct precedent for the operation-gated reflection mechanism we develop in Section~\ref{sec:reflection}.

Belkhale \emph{et al.}~\cite{belkhale2024rthactionhierarchiesusing} complemented this with RT-H, which provides an intermediate layer of language motions as an explicit action hierarchy connecting high-level task instructions with raw motor commands. This multi-level design improves data efficiency across semantically diverse multi-task datasets by forcing the model to learn shared structure at the motion level and providing an easy-to-use interface for human corrections at runtime. Because the intermediate representations are natural-language strings, a human operator can intervene by simply providing a corrected motion, which the policy then executes directly. This improves policy robustness relative to flat architectures such as RT-2~\cite{brohan2023rt2visionlanguageactionmodelstransfer}, where correcting behavior requires re-demonstrating entire trajectories.

% ------------------------------------------------------------------------------
\section{Large-Scale Deployment and Orchestration}
The field’s focus has expanded from single-task policy learning to the challenge of deploying and coordinating \gls{fm}-driven robots at scale. Ahn \emph{et al.}~\cite{ahn2024autortembodiedfoundationmodels} demonstrated, using AutoRT, that a model can autonomously coordinate entire fleets of real-world robots. AutoRT operated across more than 20 robots in four buildings over seven months. It uses a \gls{vlm} to describe each robot’s local scene and an \gls{llm} to propose situationally appropriate tasks. A major contribution of that work is the introduction of a robot constitution. It provides safety and behavioral constraints, encoded with prompts, to align autonomous task selection with human preferences.

Concurrently, Google DeepMind’s Gemini Robotics~\cite{google2026} extended this work. It integrates a general-purpose multimodal reasoning engine into real-time physical control, enabling robots to respond to unstructured human communication.

As the survey by Liu \emph{et al.}~\cite{liu2025aligningcyberspacephysical} highlights, the main challenge of the field is no longer simply aligning image representations with semantic labels. Instead, it is aligning the digital reasoning space with the physical world so that what a model understands corresponds directly to what a robot can physically execute. The reviews by Wu and Suh~\cite{wu2024stateoftheartrobotlearningmultirobot} and Kawaharazuka \emph{et al.}~\cite{Kawaharazuka_2024} converge on the conclusion that the next generation of \gls{hrc} systems will be defined by mutual adaptation. Robots should learn human preferences and adapt their behavior dynamically over time, rather than simply following instructions.

% ------------------------------------------------------------------------------
\section{Multi-Robot Coordination and Multi-Agent Planning}
Coordinating several robots toward a shared goal is a distinct problem. The system must allocate roles, sequence interdependent actions, and resolve conflicts among agents sharing a workspace. As the survey by Wu and Suh~\cite{wu2024stateoftheartrobotlearningmultirobot} documents, learning-based multi-robot collaboration has historically depended on centralized policies or hand-engineered coordination rules. Both scale poorly as the space of possible interactions grows. Language models offer an alternative, since the same reasoning that grounds one robot’s plan can, in principle, also arbitrate between several.

A single centralized planner is the simplest option, but it tends to under-specify how work is divided and to emit plans whose conflicts surface only at execution time. RoCo by Mandi \emph{et al.}~\cite{mandi2023rocodialecticmultirobotcollaboration} approaches this by giving each robot its own \gls{llm} agent. The agents then negotiate a joint plan through a dialectical exchange before any motion is run, so disagreements appear as language. This distributed formulation directly inspires the negotiation protocol we develop in Section~\ref{sec:negotiation}. Fleet-scale frameworks such as AutoRT~\cite{ahn2024autortembodiedfoundationmodels} coordinate many robots, but each acts independently on its own assignment.

% ------------------------------------------------------------------------------
\section{Synthesis and Research Gap}
One trade-off runs through all of this work. End-to-end \gls{vla} models such as PaLM-E~\cite{driess2023palmeembodiedmultimodallanguage} and RT-2~\cite{brohan2023rt2visionlanguageactionmodelstransfer} achieve impressive generalization but demand massive training data and offer limited interpretability. RT-H~\cite{belkhale2024rthactionhierarchiesusing} improves this by introducing an intermediate language-motion layer, essentially a semantic bridge between language instructions and motor actions, which enhances understandability and supports targeted human corrections without sacrificing generalization. Modular, code-generating approaches like \gls{cap}~\cite{liang2023codepolicieslanguagemodel} improve compositionality and transparency by generating code snippets that directly compose with existing perception and control APIs, but their reliability depends on the available API set and degrades when they must adapt to substantially different control interfaces. Framework-level solutions such as \gls{ros}-\gls{llm}~\cite{Mower_2026} and AutoRT~\cite{ahn2024autortembodiedfoundationmodels} provide practical scaffolding for deployment, although they rely on carefully selected action libraries and safety limitations that must be engineered per domain. The gap remains specific. No existing system fully combines the generality of \glspl{fm} with the reliability demanded by physical execution.

We set out to close this gap by separating visual understanding from execution and grounding the planner in a reusable library of basic operations, then letting that same grounding drive cooperation. The aim is genuine, autonomous cooperation between two robotic arms that follows from grounded reasoning; the following chapters detail how the framework achieves it.

